% Target ~180 words.

We propose Articulatory Trajectory Deviation (ATD), a zero-shot metric for
dysarthric speech severity that exploits the geometric structure of frozen
self-supervised speech representations. Pre-trained on 94k hours of
neurotypical speech, WavLM-Large implicitly learns a manifold of normal
articulatory execution; dysarthric speakers trace deviated trajectories
through this manifold. ATD measures the per-utterance dynamic-time-warping
distance between the speaker's WavLM trajectory and a per-word reference
trajectory built by DTW-barycentre averaging of seven control speakers,
operating in a shared 64-dimensional PCA subspace and length-normalized to
isolate execution quality from speaking rate. Without fine-tuning of any
model, ATD attains a speaker-level Spearman $\rho = 0.756$ with severity
labels on TORGO ($p = 0.0011$, $n = 15$, bootstrap CI $[0.40,\,0.92]$),
correctly ranks 7 of 8 dysarthric speakers under leave-one-speaker-out, and
reduces RMSE by 16.3\% over MFCC + SVR baselines. A controlled-perturbation
study confirms that ATD responds specifically to articulatory distortion
(formant shift, frame dropout) and not to acoustic factors (time-stretch,
additive noise). Phoneme-level path decomposition further localises
deviation to specific articulators.
